\usepackage{booktabs}
\usepackage{lastpage}
\geometry{left=2cm,right=2cm,top=3cm,bottom=2cm,headheight=3cm}
\definecolor{ceh.dark.blue}{HTML}{113ea8}
\usepackage{sectsty}
\allsectionsfont{\color{ceh.dark.blue}}
\usepackage[yyyymmdd,hhmmss]{datetime}
\usepackage{fancyhdr}
\pagestyle{fancy}
\fancyheadoffset{0cm}
\fancyhf{}
\renewcommand{\headrulewidth}{1pt}
\renewcommand{\headrule}{\hbox to\headwidth{\color{ceh.dark.blue}\leaders\hrule height \headrulewidth\hfill}}
\lhead{\includegraphics[width=4cm]{ukceh_logo_long_720x170_rgb.png}} \rhead{}
\lfoot{\thepage\ of \pageref{LastPage} |   MAVIS    |   Author   |   \today}
\fancypagestyle{plain}{\pagestyle{fancy}}
\usepackage{titling}
\setlength{\droptitle}{-2cm}
\pretitle{\begin{flushleft}\LARGE}
\posttitle{\end{flushleft}}
\preauthor{\begin{flushleft}\Large}
\postauthor{\end{flushleft}}
\predate{\begin{flushleft}}
\postdate{\end{flushleft}}
\usepackage{float}
\usepackage{longtable}
\usepackage{pdflscape}
